% MTH5000 Masters Research Project -- oral presentation (first week of November)
% Jaykumar Vinodbhai Chauhan, 34710280
%
% Filled in from report/main.tex. Every number on these slides is copied from
% the same table or script result the report itself uses (report/tables/*.tex,
% as \input or quoted directly) -- never a figure recomputed or guessed for
% the slide. Figures are not copied here: this reuses report/figures and the
% root-level fig*.png directly via \graphicspath below.
%   tectonic slides.tex

\documentclass[12pt,aspectratio=169]{beamer}

\usepackage[T1]{fontenc}
\usepackage[utf8]{inputenc}
\usepackage{mathptmx}
\usepackage{graphicx}
\usepackage{amsmath,amssymb}
\usepackage{booktabs}

\usetheme{default}
\setbeamertemplate{navigation symbols}{}
\setbeamertemplate{footline}[frame number]
\setbeamercolor{title}{fg=black}
\setbeamercolor{frametitle}{fg=black}
\setbeamercolor{structure}{fg=black}
\setbeamertemplate{itemize items}[circle]

\graphicspath{{../report/figures/}{../}}

\title{Machine Learning for Forecasting Hazardous Air Quality Days in Indian Cities}
\author{Jaykumar Vinodbhai Chauhan \\ \small Student ID: 34710280 \\ Supervisor: Dr Tianhai Tian}
\institute{School of Mathematics, Monash University \\ MTH5000 Masters Research Project}
\date{\today}

\AtBeginSection[]{
  \begin{frame}{Outline}
    \tableofcontents[currentsection]
  \end{frame}
}

\begin{document}

\begin{frame}
  \titlepage
\end{frame}

\begin{frame}{Outline}
  \tableofcontents
\end{frame}

% ============================================================================
\section{Introduction}

\begin{frame}{Motivation}
  \begin{itemize}
    \item PM2.5 in Indian cities regularly reaches hazardous levels;
      sustained exposure raises cardiopulmonary and all-cause mortality
    \item A forecast only helps if it arrives \emph{before} the hazardous
      day -- same-day warnings are of limited use
    \item Daily PM2.5 has exploitable structure: day-to-day persistence, a
      strong annual cycle, dependence on recent conditions
    \item My own motivation: growing up in India, seasonal smog was
      ordinary life, not an abstract statistic
  \end{itemize}
\end{frame}

\begin{frame}{Research Problem and Significance}
  \begin{itemize}
    \item Existing ML studies of Indian air quality report average error
      only -- not decomposed by the operating point a warning uses
    \item Good average error can still mean a poor warning system, and
      the reverse also holds
    \item The approved proposal assumed a weekly cycle alongside the
      annual one -- tested directly here, not carried over
    \item \textbf{Significance:} if the decision threshold dominates the
      model, the alarm budget is the more consequential design choice
  \end{itemize}
\end{frame}

\begin{frame}{Objectives and Research Questions}
  \textbf{Objectives}
  \begin{enumerate}
    \item Evaluate classical, ML and deep learning models 1--3 days ahead,
      under a leakage-tested protocol
    \item Compare forecast-then-threshold against direct classification
    \item Quantify skill as a warning system under a varying threshold
  \end{enumerate}
  \textbf{Research questions}
  \begin{enumerate}
    \item How accurately can PM2.5 be forecast 1--3 days ahead?
    \item Threshold a forecast, or classify the exceedance directly?
    \item What hit rate / false alarm rate trade-off does each give?
  \end{enumerate}
  \centering\textit{Q3's answer turns out to dominate Q1 and Q2.}
\end{frame}

% ============================================================================
\section{Literature Review}

\begin{frame}{Classical and Machine Learning Methods}
  \begin{itemize}
    \item \textbf{ARIMA} (Box--Jenkins): differenced series, own past
      values/errors; order by an AIC grid (Hyndman \& Khandakar, 2008);
      Fourier terms for complex seasonality (De Livera et al., 2011)
    \item \textbf{ML lineage:} ridge (Hoerl \& Kennard, 1970); random
      forest = bagging (Breiman, 1996) + trees (Breiman et al., 2017) +
      decorrelation (Breiman, 2001); gradient boosting, additive and
      stagewise (Friedman, 2001)
    \item \textbf{Competition evidence:} flexible learners need not beat
      ARIMA on short series (M3/M4: Makridakis et al., 2018/2020); an
      LSTM lost to ARIMA head-to-head (Siami-Namini et al., 2018) -- but
      gradient boosting won M5, where genuine exogenous structure existed
      (Makridakis et al., 2022)
    \item \textbf{Takeaway:} the deciding factor is exploitable
      structure, not ``ML vs.\ classical'' in the abstract -- this
      project's single series sits closer to M4 than M5
    \item A single train-test split understates variability (Tashman,
      2000; Hewamalage et al., 2023) -- the motivation for rolling-origin
      evaluation
  \end{itemize}
\end{frame}

\begin{frame}{Deep Learning: Why a GRU}
  \begin{itemize}
    \item \textbf{MLP} (Hornik et al., 1989): no sense of sequence order
      -- needs the same hand-built lag table as the regression and tree
      models
    \item \textbf{CNN} (LeCun et al., 1998): fixed receptive field,
      ill-suited to an annual cycle already handed to every model as
      Fourier terms
    \item \textbf{Transformer / TFT} (Lim et al., 2021): large parameter
      counts tuned for datasets orders of magnitude bigger than
      $\sim$3{,}000 observations
    \item \textbf{Plain RNN:} vanishing/exploding gradients over long
      sequences (Bengio et al., 1994)
    \item \textbf{LSTM} fixes this with 3 gates and a memory cell --
      more parameters than this dataset needs (Hochreiter \& Schmidhuber,
      1997)
    \item \textbf{GRU} (Cho et al., 2014): 2 gates, no memory cell, same
      gradient-flow property, smaller parameter count -- the assumptions
      that best match this project's data
  \end{itemize}
\end{frame}

\begin{frame}{Delhi Air Quality and the Research Gap}
  \begin{itemize}
    \item Delhi's emissions: traffic, industry, seasonal stubble burning
      (Guttikunda et al., 2014; Cusworth et al., 2018; Jethva et al.,
      2019)
    \item Prior Delhi PM2.5 ML studies (Masood \& Ahmad, 2020; Singh \&
      Srivastava, 2025) score accuracy only -- \textbf{none} evaluates forecasts
      as a warning system with an explicit alarm budget
    \item \textbf{Verification framework} (Jolliffe \& Stephenson, 2012):
      match measures to the decision -- CSI (Schaefer, 1990); plain
      accuracy misleads on imbalanced events (He \& Garcia, 2009; Saito
      \& Rehmsmeier, 2015)
    \item \textbf{Data quality literature} on negative readings (Jiang et
      al., 2023) and missingness mirrors the two defects found here
    \item \textbf{The gap:} evaluate against an explicit alarm budget --
      Chapter 4 shows this changes which model wins
  \end{itemize}
\end{frame}

% ============================================================================
\section{Methodology}

\begin{frame}{Data and Two Defects in the Source}
  \begin{columns}
    \column{0.58\textwidth}
    \begin{itemize}
      \item OpenAQ. Every CPCB station returns only $\sim$18 months;
        only location 8118 (AirNow, a US diplomatic post) has the full
        decade -- selected, with the limitation stated plainly
      \item \textbf{Clipped ceiling:} 1990\,\textmu g\,m$^{-3}$ repeats
        identically on 3 unrelated dates -- an instrument ceiling, not a
        reading
      \item \textbf{Wrong coverage denominator:} assumed 24 hourly
        readings, sensor was half-hourly 2016--2024. Recomputed: 1{,}085
        of 3{,}467 days are below 75\% true coverage, vs.\ 268 with no
        value at all
    \end{itemize}
    \column{0.42\textwidth}
    \centering\footnotesize
    Period: 2016--2026\\
    Usable days: 3{,}093 (86.4\%)\\
    Longest gap: 34 days\\
    Days above threshold: 920 (29.7\%)\\[4pt]
    \includegraphics[width=\textwidth]{fig2_coverage}
  \end{columns}
\end{frame}

\begin{frame}{Seasonality: Testing, Not Assuming, a Weekly Cycle}
  \begin{columns}
    \column{0.55\textwidth}
    \begin{itemize}
      \item Proposal assumed a weekly cycle alongside the annual one --
        tested directly, three independent ways
      \item \textbf{Periodogram:} dominant periods near 358, 179, 61 days
        -- the annual cycle and harmonics, nothing near 7
      \item \textbf{ACF:} lag 7 is 0.173, lag 6 is 0.179 -- on the decay
        curve, not above it
      \item \textbf{Day-of-week:} ANOVA $p=0.34$; Kruskal-Wallis $p=0.77$
      \item \textbf{No weekly cycle} -- consistent with a diplomatic
        enclave, not a roadside site. ARIMA carries annual Fourier terms
        only
    \end{itemize}
    \column{0.45\textwidth}
    \centering
    \includegraphics[width=0.95\textwidth]{fig3_weekly}
  \end{columns}
\end{frame}

\begin{frame}{Features and Weather Covariates}
  \begin{itemize}
    \item One row per origin $t$, using only data up to day $t$: lags,
      rolling mean/SD/max, anomalies, exceedance frequency, calendar and
      annual Fourier terms, missingness flags
    \item Missing values are never imputed, except for the GRU, shown
      next
    \item \textbf{Weather} (\texttt{meteostat}): temperature, wind speed,
      pressure, precipitation, each with a lag and a 3-day rolling mean
    \item \textbf{Third data defect:} the weather archive ends 202 days
      before the PM2.5 series -- a hard cutoff. Precipitation is only
      54\% complete and excluded from the GRU
  \end{itemize}
\end{frame}

\begin{frame}{Classical Benchmark and Machine Learning Regressors}
  \textbf{ARIMA(1,1,1)} with four pairs of annual Fourier terms, order by
  AIC (1892.5):
  \[
  (1-\phi_1 B)(1-B)y_t = c + (1+\theta_1 B)\varepsilon_t +
  \sum_{k=1}^{4}\big[\alpha_k\sin(\tfrac{2\pi kt}{365.25}) +
  \beta_k\cos(\tfrac{2\pi kt}{365.25})\big]
  \]
  fitted on $\log(\text{PM2.5})$; $\hat\phi_1=0.581$, $\hat\theta_1=-0.987$.

  \textbf{Three ML regressors}, each fitted to level and change from
  origin:
  \begin{itemize}
    \item \textbf{Ridge:} closed-form $\hat\beta=(X^\top X+\lambda I)^{-1}X^\top y$
    \item \textbf{Random forest:} bagged trees + random feature subset,
      which decorrelates trees and lowers the variance floor
    \item \textbf{Gradient boosting:} additive, tree per round corrects
      the residual; accepts missing values natively
  \end{itemize}
\end{frame}

\begin{frame}{The Gated Recurrent Unit}
  \begin{columns}
    \column{0.5\textwidth}
    \footnotesize
    \begin{align*}
      z_t &= \sigma(W_z x_t + U_z h_{t-1} + b_z)\\
      r_t &= \sigma(W_r x_t + U_r h_{t-1} + b_r)\\
      \tilde h_t &= \tanh(W_h x_t + U_h(r_t\odot h_{t-1}) + b_h)\\
      h_t &= (1-z_t)\odot h_{t-1} + z_t\odot \tilde h_t
    \end{align*}
    \column{0.5\textwidth}
    \footnotesize
    \begin{itemize}
      \item $H=24$ hidden units, 30-day window
      \item 6 channels (10 with weather): log-concentration, missingness
        flag, Fourier pairs ($K=2$ vs.\ ARIMA's $K=4$)
      \item Gaps forward-filled, missingness flagged
      \item Adam, MAE loss, early stopping
      \item \textbf{10 seeds} -- not bit-reproducible at a fixed seed;
        every result is a range
    \end{itemize}
  \end{columns}
\end{frame}

\begin{frame}{Rolling-Origin Evaluation and Warning-System Measures}
  \begin{itemize}
    \item \textbf{2{,}303 origins}, Nov 2019--Aug 2026, refit every 90
      days (180 for the GRU). At origin $t$, only data up to day $t$ is
      used -- enforced by construction, verified by a corruption test
    \item \textbf{Accuracy:} rMAE against persistence (1.000 =
      persistence) and MASE, per horizon
    \item \textbf{Warning system:} hit rate, false alarm rate, precision,
      CSI, alarms per year -- plain accuracy is never reported
    \item \textbf{Diebold-Mariano test} (HLN small-sample correction):
      GRU vs.\ ARIMA per seed, and ARIMA vs.\ each ML regressor
  \end{itemize}
\end{frame}

% ============================================================================
\section{Results}

\begin{frame}{Forecast Accuracy (RQ1)}
  \centering\footnotesize
  \begin{tabular}{lccc}
    \toprule
    Model & $h=1$ & $h=2$ & $h=3$ \\
    \midrule
    GRU, range over 10 seeds & 0.910--0.939 & 0.878--0.906 & 0.831--0.879 \\
    ARIMA(1,1,1) + Fourier & \textbf{0.940} & \textbf{0.883} & \textbf{0.832} \\
    Best ML regressor (GBM, change) & 0.983 & 0.937 & 0.878 \\
    Persistence & 1.000 & 1.000 & 1.000 \\
    \bottomrule
  \end{tabular}

  \includegraphics[width=0.46\textwidth]{fig6_accuracy}

  \raggedright\small
  GRU wins clearly at $h=1$; ARIMA wins or ties beyond it ($h=3$: 0.831
  vs.\ 0.832). ARIMA beats every ML regressor at every horizon: 15 of 15
  DM pairings favour it, 14 at $p<0.05$.
\end{frame}

\begin{frame}{Does Weather Change the Answer?}
  \begin{itemize}
    \item Gradient boosting improves at every horizon (by 8.2 / 2.0 / 2.9
      points of rMAE); the random forest worsens 5.3 points at $h=3$ --
      no single story for every model
    \item \textbf{GRU + weather:} wins clearly at $h=1$ (7/10 seeds now
      significant, up from 4/10); $h=2$ becomes a genuine 3-way split;
      $h=3$ is more often a loss than a win (5 of 10)
    \item Doubling the seed count (5$\to$10) is what surfaced this --
      under five seeds alone, $h=2$ and $h=3$ had looked like wins
    \item \textbf{Carried through the threshold:} at most 1.2 points of
      extra hit rate at any budget -- a materially more accurate
      forecaster still does not escape the alarm-budget ceiling
  \end{itemize}
\end{frame}

\begin{frame}{Two Arms of the Same Question (RQ2)}
  \begin{columns}
    \column{0.56\textwidth}
    \begin{itemize}
      \item Forecast-then-threshold vs.\ direct classification, matched
        by base algorithm
      \item 36 comparisons (3 pairs $\times$ 3 horizons $\times$ 4
        budgets): \textbf{threshold wins 21, classifier 11, 4 ties}
      \item Median gap 0.011 hit rate, max 0.056 -- real but modest
      \item Edge narrows as budget loosens: 8/9 at budget 40, but the
        classifier overtakes it (6/9) at budget 100
      \item Preferred anyway: a forecast can be rethresholded for any
        future policy; a classifier is welded to its training threshold
    \end{itemize}
    \column{0.44\textwidth}
    \includegraphics[width=\textwidth]{fig7_arms}
  \end{columns}
\end{frame}

\begin{frame}{The Decision Threshold (RQ3) -- the Central Result}
  \centering\footnotesize
  \begin{tabular}{lrrrr}
    \toprule
    Alarms per year & 40 & 60 & 80 & 100 \\
    \midrule
    Best hit rate, $h=1$ & 0.379 & 0.567 & 0.735 & 0.881 \\
    Best hit rate, $h=2$ & 0.366 & 0.538 & 0.700 & 0.840 \\
    Best hit rate, $h=3$ & 0.359 & 0.533 & 0.685 & 0.827 \\
    \bottomrule
  \end{tabular}

  \includegraphics[width=0.38\textwidth]{fig5_warning}

  \raggedright\small
  Budget 40$\to$100 moves $h=1$ hit rate 0.379$\to$0.881: \textbf{50
  points}. Switching the model instead moves it \textbf{$\sim$5 points}.
  At $h=3$, budget 40, plain persistence is the best model in the study.
\end{frame}

% ============================================================================
\section{Discussion}

\begin{frame}{The Operational Question Dominates the Modelling Question}
  \begin{itemize}
    \item Alarm budget moves the hit rate $\sim$50 points; switching the
      model at fixed budget moves it $\sim$5
    \item At $h=3$, loosest budget: climatology alone (0.816) comes
      within a hundredth of the best fitted model (0.827)
    \item \textbf{Not} an argument against modelling -- ARIMA and the GRU
      both beat persistence, and a better forecast narrows the useful
      range of thresholds
    \item The leverage sits with the policy instrument, the alarm budget,
      not the model behind it. Recommendation: fix the budget with the
      health authority first; treat the model as a secondary refinement
  \end{itemize}
\end{frame}

\begin{frame}{How Large Is a Difference Between Models?}
  \begin{itemize}
    \item \textbf{Noise floor:} identical pipeline, two machines, moves
      tree models up to 0.004; the GRU moves 0.028--0.048 across seeds --
      almost as much when the \emph{same} seed is simply rerun
    \item Several table gaps sit at or below that floor (e.g.\ 0.878
      vs.\ 0.879 at $h=3$) and cannot be confidently ordered
    \item \textbf{Diebold-Mariano is more conservative still:} only 4/10
      GRU seeds reach $p<0.05$ at $h=1$ (the noise floor credits 8/10);
      0/10 at $h=2$; 1/10 at $h=3$
    \item Those seed counts are themselves small-sample proportions: a
      95\% Wilson interval on ``4 of 10'' at $h=1$ runs 0.17 to 0.69
  \end{itemize}
\end{frame}

\begin{frame}{Limitations I: Scope and Sensitivity}
  \begin{itemize}
    \item \textbf{Single station, single city:} an AirNow diplomatic-post
      monitor, not the CPCB network; Delhi only, chosen for data history
    \item \textbf{Cleaning threshold} (1000\,\textmu g\,m$^{-3}$) checked
      directly: restoring capped days barely moves ARIMA (0.940$\to$0.935
      at $h=1$) but hurts the ML suite more (ridge: 0.997$\to$1.118) --
      the central finding \emph{strengthens}
    \item Tree hyperparameters from a small grid, not an exhaustive
      search -- a lower bound on what those models could achieve
  \end{itemize}
\end{frame}

\begin{frame}{Limitations II: The GRU and Reproducibility}
  \begin{itemize}
    \item GRU architecture ($H=24$, 30-day window) fixed by inspection --
      a bounded check found a 45-day window scored 2.2 points better at
      $h=3$, untested at full ten-seed scale
    \item \textbf{Fourier ablation:} giving the GRU ARIMA's $K=4$
      harmonics helps at every horizon, substantially at $h=3$ (3.96
      points) -- the network already gets a seasonal signal, just a
      coarser one
    \item GRU training is \textbf{not bit-reproducible} at a fixed seed
      -- why the seed count was doubled from 5 to 10
    \item Weather is observed only, never forecast; boundary layer
      height is unavailable without a reanalysis product
  \end{itemize}
\end{frame}

% ============================================================================
\section{Conclusion}

\begin{frame}{Summary and Main Conclusions}
  \begin{itemize}
    \item Forecast hazardous PM2.5 1--3 days ahead at one Delhi monitor;
      two real data defects corrected; the weekly-cycle assumption tested
      and rejected
    \item Spread across models at fixed budget: $\sim$5 points. Spread
      across budgets, model fixed: $\sim$50 points
    \item ARIMA beats every ML regressor at every horizon (DM-backed);
      the GRU is the one model that wins, only at $h=1$, reported as a
      range over ten seeds
    \item Weather narrows $h=1$ further but leaves $h=2$/$h=3$ a tie or
      more-often-loss; its gain buys little once thresholded
    \item Forecast-then-threshold beats direct classification 21 of 36 --
      real but modest, preferred also for rethresholdability
    \item Trying several methods, not defaulting to ARIMA, is what
      surfaced the $h=1$ GRU win and made the negative result credible
  \end{itemize}
\end{frame}

\begin{frame}{Future Work}
  \begin{enumerate}
    \item A genuine weather \emph{forecast}, not the conditions already
      observed at the forecast origin
    \item Boundary layer height via a reanalysis product (e.g.\ ERA5)
    \item Refit the GRU with four Fourier harmonics under the full
      ten-seed rolling-origin evaluation
    \item A second, independent station -- ideally CPCB -- once its
      record lengthens
  \end{enumerate}
  \vspace{6pt}
  \centering
  \textit{Design the warning system around its alarm budget; treat the
  forecasting model as a secondary refinement.}
\end{frame}

\begin{frame}
  \centering
  \Large Questions?
\end{frame}

\end{document}
